\newif\ifglobal

%%%%%%%%%%%%%%%%%%%%%%%
%%% Compile to view %%%
%%%%%%%%%%%%%%%%%%%%%%%

%%%%%%%%%%%%%%%%%%%%%%%%%%%%%%%%%%%%%%%%%%%%%%%%%%%
%%% Readme:                                     %%%
%%% https://www.overleaf.com/read/bwdjvfjksvjy/ %%%
%%%%%%%%%%%%%%%%%%%%%%%%%%%%%%%%%%%%%%%%%%%%%%%%%%%

% >>> M?: marks a module setting number or important notice

% >>> M4: Set {./relative-path} to external files for this module <<<
\def \modulefiles {./35-RNG}
% To add an external file, use:
% (Replace '---' with actual filename)
% '\input{\modulefiles/tables/---.tex}' for tables
% '\includegraphics{\modulefiles/figures/---}' for figures
% '\subfile{\modulefiles/---}' for register files

% >>> M1: This module is in global project? <<<
% 'YES' - uncomment; 'NO' - comment out
% \globaltrue

\ifglobal
    \documentclass[main\_\_EN.tex]{subfiles}

\else
    \documentclass[openany, 10pt]{book}

    % >>> M2: In ./chip-spec-settings/preamble-trm-module, also find
    %         settings for visibility of todo notes, watermark <<<
    \usepackage{preamble-trm-module}

    % >>> M3: Temporary labels in English,
    %         you can add your own labels there <<<
    \myexternaldocument{temp-labels-en}

\fi %\ifglobal


\setcounter{secnumdepth}{4}




% >>> M5: If this module needs any updates in
% './chip-spec-settings/preamble-trm-module.sty'
% (include additional package, etc.), contact your module PM
% or create the file './chip-spec-settings/changelog.tex' make the
% necessary updates yourself and document them in this file <<<

\raggedright


\begin{document}


% Sets language into which pre-defined bits of text are translated
\selectlanguage{english}

% Do not delete this block
\ifglobal
% Add the following front matter if
% this module is outside of global project
\else
    \InputIfFileExists{readme}{}{}
    {\let\clearpage\relax \listoftodos}
    {\let\clearpage\relax \listofdonetodos}
    \tableofcontents
    %\listoftables
    %\listoffigures
\fi


%%%%%%%%%%%%%%%%%%%%%%%%%
%%% TRM Module Begins %%%
%%%%%%%%%%%%%%%%%%%%%%%%%

\hypertarget{rng}{}
\chapter{Random Number Generator (RNG)}
\label{mod:rng}


\section{Overview}\label{sec:rng-overview}

The \chipname{} has a true random number generator that can generate two types of random numbers: CRC random numbers and NIST random numbers. Both types of random numbers can be used as the basis for cryptographic operations, and both have passed Dieharder tests.

%%%%%%%%%%%%%%%%%%%
%%% Terminology %%%
%%%%%%%%%%%%%%%%%%%

\section{Terminology}
\label{sec:rng-term}

To better illustrate the functionality of RNG, the following terms are used in this chapter.

\begin{longtable}[c]{ R{4.5cm} p{11.5cm} }
    \textbf{CRC random numbers} & Random numbers generated by the true random number generator and readable by software from register \hyperref[regdesc:TRNGCRCSYNCDATAREG]{TRNG\_CRC\_SYNC\_DATA\_REG}.\\
    \textbf{NIST random numbers} & Random numbers generated by the true random number generator for cryptographic modules in the chip security subsystem. They are checked by the on-chip health-test circuit that complies with NIST SP 800-90B.\\
    \textbf{NIST SP 800-90B} & National Institute of Standards and Technology Special Publication 800-90B. This publication defines the health tests used by RNG: Startup Test, Repetition Test, and Adaptive Test.\\
\end{longtable}

%%%%%%%%%%%%%%%%%%%%
%%% Feature List %%%
%%%%%%%%%%%%%%%%%%%%

\section{Feature List}
\label{sec:rng-features}

\begin{itemize}
    \item Thermal noise, asynchronous clocks, and a ring oscillator as noise sources
    \item CRC random numbers readable by software
    \item NIST random numbers for cryptographic modules in the chip security subsystem
    \item On-chip health-test circuit compliant with NIST SP 800-90B (see Section \ref{sec:rng-nist-rng} \textit{\nameref{sec:rng-nist-rng}})
\end{itemize}

%%%%%%%%%%%%%%%%%%%%%%%%%%%%%%
%%% Functional Description %%%
%%%%%%%%%%%%%%%%%%%%%%%%%%%%%%

\section{Functional Description}
\label{sec:rng-func-descr}

Both CRC random numbers and NIST random numbers can use \textbf{thermal noise}, \textbf{asynchronous clocks}, and \textbf{a ring oscillator} as noise inputs. In addition, the NIST path supports reusing CRC\_DATA as an input.

\begin{itemize}
    \item \textbf{Thermal noise} can come from SAR\_ADC or RF\_ADC. When SAR\_ADC is running, SAR\_ADC generates a bit stream that is used as a random seed for the random number generator through XOR logic.
    \item RC\_FAST\_CLK is an \textbf{asynchronous clock source} that introduces circuit metastability. This metastability can also be used as a random seed. For RC\_FAST\_CLK, see the \hyperref[others:rng-note-clock-seed]{Note} below.
    \item The count value generated by the \textbf{ring oscillator} (BUF\_CHAIN), implemented using BUF\_CELL, is used as a random seed for the random number generator.
\end{itemize}
The noise-source structure is shown in Figure \ref{fig:rng-noise-source} \textit{\nameref{fig:rng-noise-source}}.
\vspace{2em}

\begin{figure}[H]
    \centering
    \includegraphics[width=1.0\textwidth]{\modulefiles/figures/35-RNG-noise-source.png}
    \caption{Noise Source}
    \label{fig:rng-noise-source}
\end{figure}

Software can read CRC random numbers from register \hyperref[regdesc:TRNGCRCSYNCDATAREG]{TRNG\_CRC\_SYNC\_DATA\_REG}. NIST random numbers have higher quality than CRC random numbers. NIST random numbers can only be requested through cryptographic modules in the chip security subsystem (for example, Key Manager), and cannot be read directly through the software path.

When SAR\_ADC is enabled, the random number generator gets 1-bit entropy in each RC\_FAST\_CLK cycle (from the internal RC oscillator, see Chapter \ref{mod:resclk} \textit{\nameref{mod:resclk}}). Therefore, to maximize entropy, Espressif recommends reading register \hyperref[regdesc:TRNGCRCSYNCDATAREG]{TRNG\_CRC\_SYNC\_DATA\_REG} at a rate no higher than 1 MHz. Reading at a higher rate may reduce randomness quality.



\section{Programming Procedures}\label{sec:rng-programm-descr}

\subsection{CRC RNG}

When using the random number generator in \chipname{}, at SoC level, complete the following steps before any TRNG operation:

\begin{enumerate}
    \item Configure \hyperref[fielddesc:LPPERICLKRSTLPRNGCKEN]{LP\_PERICLKRST\_LP\_RNG\_CLK\_EN} in register \hyperref[regdesc:LPPERICLKRSTRNGCTRLREG]{LP\_PERICLKRST\_RNG\_CTRL\_REG} to 1 to turn on the RNG clock and the RC\_FAST\_CLK which is used as an entropy source.
    \item Configure \hyperref[fielddesc:LPPERICLKRSTLPRNGRSTEN]{LP\_PERICLKRST\_LP\_RNG\_RST\_EN} to 1, then configure it to 0, to generate a reset pulse for the RNG IP.
\end{enumerate}

Then enable BUF\_CHAIN. Otherwise, pseudo-random numbers may be generated. SAR\_ADC and RC\_FAST\_CLK may also be enabled as additional noise sources.
\begin{itemize}
    \item Configure \hyperref[fielddesc:TRNGSAMPLEENABLE]{TRNG\_SAMPLE\_ENABLE} in register \hyperref[regdesc:TRNGCONFREG]{TRNG\_CONF\_REG} to enable BUF\_CHAIN.
    \item For SAR\_ADC, see Chapter \ref{mod:adc} \textit{\nameref{mod:adc}}.
    \item Configure \hyperref[fielddesc:PMUHPSLEEPXPDFOSCCLK]{PMU\_HP\_SLEEP\_XPD\_FOSC\_CLK} in register \hyperref[regdesc:PMUHPSLEEPLPCKPOWERREG]{PMU\_HP\_SLEEP\_LP\_CK\_POWER\_REG} to enable RC\_FAST\_CLK. For RC\_FAST\_CLK, see the \hyperref[others:rng-note-clock-seed]{Note} below.
\end{itemize}

Then set \hyperref[fielddesc:TRNGNOISECRCEN]{TRNG\_NOISE\_CRC\_EN} to 1 to enable the CRC random number generation.

\begin{tiplisting}\label{others:rng-note-clock-seed}
    RC\_FAST\_CLK can improve RNG entropy. However, to ensure sufficient entropy, Espressif still recommends keeping BUF\_CHAIN enabled when RC\_FAST\_CLK is used as a noise source.
\end{tiplisting}

When software uses the random number generator, read \hyperref[regdesc:TRNGCRCSYNCDATAREG]{TRNG\_CRC\_SYNC\_DATA\_REG} multiple times until enough random numbers are obtained. During reading, keep the read rate within the limit described in Section \ref{sec:rng-func-descr} \textit{\nameref{sec:rng-func-descr}}.

\subsection{NIST RNG}\label{sec:rng-nist-rng}

For NIST random numbers, select noise sources with \hyperref[fielddesc:TRNGNOISESOURCESEL]{TRNG\_NOISE\_SOURCE\_SEL} and \hyperref[fielddesc:TRNGNOISEPOSSEL]{TRNG\_NOISE\_POS\_SEL} in register \hyperref[regdesc:TRNGCONFREG]{TRNG\_CONF\_REG}:
\begin{itemize}
    \item Reuse CRC random numbers as the NIST health-test input random data: set \hyperref[fielddesc:TRNGNOISESOURCESEL]{TRNG\_NOISE\_SOURCE\_SEL}[4] = 1. If this is selected, all noise sources listed below are ignored.
    \item Enable SAR\_ADC: set \hyperref[fielddesc:TRNGNOISESOURCESEL]{TRNG\_NOISE\_SOURCE\_SEL}[3] = 1.
    \item Enable RF\_ADC: set \hyperref[fielddesc:TRNGNOISESOURCESEL]{TRNG\_NOISE\_SOURCE\_SEL}[2] = 1.
    \item Enable RC\_FAST\_CLK: set \hyperref[fielddesc:TRNGNOISESOURCESEL]{TRNG\_NOISE\_SOURCE\_SEL}[1] = 1.
    \item Enable BUF\_CHAIN: set \hyperref[fielddesc:TRNGNOISESOURCESEL]{TRNG\_NOISE\_SOURCE\_SEL}[0] = 1.
\end{itemize}

Except for the CRC\_DATA path, the noise sources above support selecting multiple input channels at the same time. To ensure output quality, if only BUF\_CHAIN or RC\_FAST\_CLK is selected, \hyperref[fielddesc:TRNGHEALTHTESTERRORCODE]{TRNG\_HEALTH\_TEST\_ERROR\_CODE} records error code 1 and clears the NIST random numbers stored in the RNG module, so that non-compliant random numbers are not output. After this error, set \hyperref[fielddesc:LPPERICLKRSTLPRNGRSTEN]{LP\_PERICLKRST\_LP\_RNG\_RST\_EN} to 1 to clear all current RNG state, and then repeat the steps that start the NIST RNG path.

For \hyperref[fielddesc:TRNGNOISEPOSSEL]{TRNG\_NOISE\_POS\_SEL}, this field uses one-hot encoding. Unlike \hyperref[fielddesc:TRNGNOISESOURCESEL]{TRNG\_NOISE\_SOURCE\_SEL}, only one bit can be set. Ensure that \hyperref[fielddesc:TRNGNOISEPOSSEL]{TRNG\_NOISE\_POS\_SEL} matches one enabled source in \hyperref[fielddesc:TRNGNOISESOURCESEL]{TRNG\_NOISE\_SOURCE\_SEL}; otherwise the RNG cannot output random numbers correctly.

To ensure NIST random number output complies with NIST SP 800-90B, RNG has three mandatory tests defined by NIST SP 800-90B: Startup Test, Repetition Test, and Adaptive Test. For \hyperref[fielddesc:TRNGREPETITIONVALUEC]{TRNG\_REPETITION\_VALUE\_C} and \hyperref[fielddesc:TRNGADAPTIVEVALUEC]{TRNG\_ADAPTIVE\_VALUE\_C} values, refer to the NIST SP 800-90B manual for recommended values based on random-number length and bit width.
\begin{itemize}
    \item Startup Test: after RNG startup, this test checks entropy-source output data sampled for a period of time to avoid output distortion caused by unstable startup behavior. Start this test by setting \hyperref[fielddesc:TRNGSTARTUPTESTSTART]{TRNG\_STARTUP\_TEST\_START} = 1. The test data length is 512 + \hyperref[fielddesc:TRNGSTARTUPTESTLIMIT]{TRNG\_STARTUP\_TEST\_LIMIT}.
    \item Repetition Test: this test checks whether output random numbers are stuck at a fixed value. It judges quality by counting how many identical values appear consecutively. Set the threshold by configuring \hyperref[fielddesc:TRNGREPETITIONVALUEC]{TRNG\_REPETITION\_VALUE\_C}. The maximum supported value is 511. A lower configured value makes the test stricter.
    \item Adaptive Test: this test checks whether a specific random value appears too frequently within any sliding window of 512 generated random numbers. Configure the limit through \hyperref[fielddesc:TRNGADAPTIVEVALUEC]{TRNG\_ADAPTIVE\_VALUE\_C}.
\end{itemize}

Set \hyperref[fielddesc:TRNGHEALTHTESTBYPASS]{TRNG\_HEALTH\_TEST\_BYPASS} to 1 to bypass the health test above only in test mode. Otherwise \hyperref[fielddesc:TRNGHEALTHTESTBYPASS]{TRNG\_HEALTH\_TEST\_BYPASS} should always be set to 0.

To read NIST random numbers from software, configure the NIST-related settings and start Startup Test, then configure \hyperref[fielddesc:TRNGSWRANDOMREQ]{TRNG\_SW\_RANDOM\_REQ}. Wait until \hyperref[fielddesc:TRNGSWDATAVALID]{TRNG\_SW\_DATA\_VALID} is 1, and then read \hyperref[fielddesc:TRNGSWRANDOMDATA]{TRNG\_SW\_RANDOM\_DATA}.

%%%%%%%%%%%%%%%%%%%%%%%%
%%% Register Summary %%%
%%%%%%%%%%%%%%%%%%%%%%%%

\newpage
\hypertarget{rng-reg-summ}{}
\section{Register Summary}
\label{sec:rng-reg-summ}




The abbreviations given in Column \textbf{Access} are explained in Section \hyperref[glossary-access-types]{\textit{Access Types for Registers}}.

\subfile{\modulefiles/35-RNG-reg-summ__EN}



%%%%%%%%%%%%%%%%%
%%% Registers %%%
%%%%%%%%%%%%%%%%%

\section{Registers}
\label{sec:rng-reg}

For how to program reserved fields, please refer to Section \hyperref[programming-reserved-field]{Programming Reserved Register Field}.

\subfile{\modulefiles/35-RNG-reg__EN}

\end{document}
